%% kapa: exam=1 round=27 year=2016 period=1 form=A subject=real_estate subject_ko=부동산학원론 seq=29 number=109
%% source: Q-Net 감정평가사 기출문제 (cst00309) artlSeq=5169402 file=extracted/1008/5169402/제27회 감정평가사 1차 문제/제27회 감정평가사 1교시문제(A형).pdf
%% conversion: kordoc 4.24.51 -> tools/gen1.py (auto); review=proofread-sonnet; flags=none
\begin{question}{109}
부동산투자회사법령상 부동산투자회사에 관한 설명으로 옳은 것은?
\begin{choices}
  \item 영업인가를 받은 날부터 6개월이 지난 자기관리 부동산투자회사의 자본금은 70억원 이상이 되어야 한다.
  \item 위탁관리 부동산투자회사 및 기업구조조정 부동산투자회사의 설립 자본금은 10억원 이상으로 한다.
  \item 자기관리 부동산투자회사의 설립 자본금은 3억원 이상으로 한다.
  \item 영업인가를 받은 날부터 6개월이 지난 위탁관리 부동산투자회사 및 기업구조조정 부동산투자회사의 자본금은 100억원 이상이 되어야 한다.
  \item 부동산투자회사는 부동산 등 자산의 운용에 관하여 회계처리를 할 때에는 국토 교통부가 정하는 회계처리기준에 따라야 한다.
\end{choices}
\end{question}
\answer{1}
